\documentclass[11pt,a4paper]{article}
\usepackage[T1]{fontenc}
\usepackage{isabelle,isabellesym}

% this should be the last package used
\usepackage{pdfsetup}

% urls in roman style, theory text in math-similar italics
\urlstyle{rm}
\isabellestyle{literal}


\begin{document}

\title{Earley With Time}
\author{Alexander Haberl and Tobias Nipkow}
\maketitle

\begin{abstract}
This is a verified development of both an executable recognizer and parser
implementing Earley's algorithm. It is accompanied by verified complexity
analyses parameterized by the time it takes to access the nth element of a list.
The asymptotic complexities of the recognizer and parser are the same:
cubic, if nth is constant-time, or quartic, if nth is linear-time.
Therefore our purely functional implementation has quartic complexity. 

Restrictions: the grammar must be epsilon-free and the parser returns only
a single parse tree.
\end{abstract}

\newpage
\tableofcontents

\newpage
\section{Introduction}

Earley's algorithm \cite{Earley-CACM} is a parsing algorithm for arbitrary context-free
languages. The recognizer only solves the word problem and does not return a parse tree.
Jones~\cite{Jones} was the first to present a rigorous correctness proof of the recognizer.
Rau and Nipkow~\cite{RauN-ITP24} verified an executable Earley parser that follows the work
of Scott~\cite{Scott-ENTCS08} and returns a shared packed parse forest (SPPF).
They verify soundness but not completeness and argue that their implementation has quartic complexity.
Nipkow and Rau~\cite{NipkowR-CBJ24} also develop separately the beginning
of a streamlined, stepwise development of Earley's recognizer.

The current AFP entry includes the theory by Nipkow and Rau
and completes the refinement to both an executable recognizer and parser.
Moreover, it contains verified complexity analyses
parameterized by the time it takes to access the nth element of a list.
The analysis is carried out precisely, with all the constants.
It employs the time function infrastructure described elsewhere \cite{Nipkow-ACMbook}.

The asymptotic complexities of the recognizer and parser are the same:
cubic, if nth is constant-time, or quartic, if nth is linear-time.
Therefore our purely functional implementation has quartic complexity,
confirming the informal argument by Rau and Nipkow~\cite{RauN-ITP24}.

In a last step, the precise analysis is abstracted with the help of a Big O notation
in order to present the asymptotic complexities concisely.

% include generated text of all theories
\newpage
\input{session}

\newpage
\bibliographystyle{abbrv}
\bibliography{root}
\addcontentsline{toc}{section}{References}

\end{document}
